\documentclass[10pt,a4paper]{amsart}
\usepackage[margin=19mm]{geometry}
\usepackage{amsmath,amssymb,mathtools}
\usepackage[scr=rsfs,cal=euler]{mathalfa}
\usepackage{xfrac}
\usepackage[unicode,hidelinks,bookmarks=false]{hyperref}
\newcommand{\ie}{i.e.\ }
\newcommand{\cf}{cf.\ }
\newcommand{\resp}{respectively\ }
\newcommand{\Subsection}{\subsection}
\providecommand{\nobreakdash}{\nobreak\hbox{-}}
\setlength{\emergencystretch}{2em}
\multlinegap0pt
\usepackage{amssymb}
\usepackage{float}
\usepackage{mathtools}
\usepackage{tikz}
\newtheorem{lemma}{Lemma}
\title{HIERARCHICAL LOCAL NULL CONTROLLABILITY FOR A NON-DEGENERATE QUASI-LINEAR PARABOLIC SYSTEM}
\author{Cristian Loli, George J. Bautista, Juan Limaco, Rafael Lobosco, Dany Nina-Huaman,, Luis P. Yapu}
\date{Source: arXiv:2608.24379v2 (CC BY 4.0)}
\begin{document}
\maketitle

\noindent\emph{Source and licence.} HIERARCHICAL LOCAL NULL CONTROLLABILITY FOR A NON-DEGENERATE QUASI-LINEAR PARABOLIC SYSTEM. Version arXiv:2608.24379v2 (CC BY 4.0), distributed by arXiv under the Creative Commons Attribution 4.0 International licence, http://creativecommons.org/licenses/by/4.0/, as recorded in the arXiv licence metadata for this exact version and shown on the abstract page https://arxiv.org/abs/2608.24379v2. Full licence notice: \texttt{licenses/arxiv-2608.24379-CC-BY-4.0.txt}.\par\smallskip

\noindent\textit{Editorial scope.} Counted unit: the article's lemma lema2 (source0.tex lines 1735--1737). The article's complete original proof of it is delivered with it (source0.tex lines 1738--1990, one \texttt{proof} environment of 253 lines). No mathematics is written, repaired or re-derived: the statement and the proof are the article's own lines, in the article's own notation, and no retained line is substituted or re-flowed.  Retained uncounted as the source-local context the proof needs: lines 180--208; lines 752--755.  Nothing was omitted from any retained span, so the package carries no omission record.  No retained line contains a citation, so the wrapper carries no bibliography and no bibliography entry.  No retained line contains a figure environment, an image-inclusion command or drawing code, so no image is delivered and the package declares no asset dependency.  Editorial formatting only, in addition: an AMS \texttt{amsart} preamble that restates the article's own declarations and macros used by the retained lines, namely the environments \texttt{lemma} on separate counters, together with the standard packages those lines need.  The retained environments print in this document as Lemma 1 in order of appearance, whereas the article prints the same items with the numbering of its own full document.  The complete native source of the article as distributed in the arXiv e-print for this exact version is bundled unchanged as \texttt{sources/208468/source0.tex}.
 We first introduce a few considerations. Let us  suppose that $\omega,\omega_1,\omega_2\subset I$ be nonempty open subsets such that
\[
\omega_i\cap\omega=\varnothing,
\qquad i=1,2.
\]
We assume that the initial data satisfy
\[
(f,g)\in [L^2(I)]^2.
\]
We impose the following assumptions on the coefficients and nonlinearities.
The diffusion coefficient satisfies
    \begin{equation} \label{condiciones_a4}
    \begin{cases}
        a\in C^3(I),\quad 0 < M_0 \leq a(r) \leq M_1,\\
        |a^{(j)}(r)|\leq M,\quad j=1,2,3,
    \end{cases}
    \end{equation}
    for suitable positive constants $M_0,M_1,M$. Moreover,
    \begin{equation}
        \begin{cases} \label{condiciones_a4f}
 		 F_i\in C^2(\mathbb{R}), \quad F_i(0)=0,\\
             |{F_{i}^{(j)}(r)}|\leq M, \text{ for all } r \in \mathbb{R}, \, i=1,2,\quad j=1,2.
 	      \end{cases}
    \end{equation}
   Finally, the coupling coefficient $c\in L^\infty(Q)$ is assumed to satisfy
   \begin{equation}\label{condiciones_a4c(x,t)}
    0<c_0<c(x,t), \text{ in } (\omega_d\cap\omega)\times(0,T),
    \end{equation}
    where $c_0>0$ is a suitable constant and  the set $\omega_d$ is defined in Section 4.

\begin{equation}\label{eq:compara_rhos3}
	\hat\rho_1 \leq C\hat\rho_0 \leq  C\rho_{1}\leq C\rho_{0}\leq C\tilde \rho, \qquad %\hat{\rho}^{2}=\rho_{1}\rho_{0} \qquad
    \text{and} \qquad \tilde \rho \leq C \hat\rho_1^2.   
\end{equation}

\begin{lemma}\label{lema2}
    The mapping $A : Y \rightarrow Z$ is continuously differentiable.
\end{lemma}

\begin{proof}
    First, we prove that $A$ is Gâteaux derivative at any $(y_1,y_2,p_i^j,h) \in Y$ and compute the
    $\textit{G-derivative}$ $A^{\prime}(y_1,y_2,p_i^j,h)$.
    Consider the linear mapping $D A: Y \to Z$ given in components by
    $$
        D A(y_1,y_2,p_i^j,h) = (D A^1_0,D A^2_0, D A^1_1, D A^2_1, D A^1_2, D A^2_2, D A^1_3, D A^2_3),
    $$
    where, for $i=1,2$,
    \begin{equation*}
    	%\left\{\begin{split}
        \begin{cases}
		      D A^1_0(\bar y_1,\bar y_2,\bar p_i^j,\bar h) = &  \,
              \bar y_{1,t} - \frac{\partial}{\partial x}(a'(y_{1})\bar y_{1} y_{1,x} +a(y_1) \bar y_{1,x}) \\
            &+  D_1 F_1(y_1,y_2)\bar y_1 + D_2 F_1(y_1,y_2)\bar y_2 - \bar h \chi_{\omega} \\
            &+ \frac{1}{\mu_1} \rho_*^{-2} \bar p^1_1 \chi_{\omega_{1}} + \frac{1}{\mu_2} \rho_*^{-2} \bar p^2_1 \chi_{\omega_{2}}, \\
		      D A^2_0(\bar y_1,\bar y_2,\bar p_i^j,\bar h) =
            &  \, \bar y_{2,t} - (a'(y_{2})\bar y_{2}y_{2,x} + a(y_{2}) \bar y_{2,x})_x \\
            &+ D_1 F_2(y_1,y_2)\bar y_1 + D_2 F_2(y_1,y_2)\bar y_2, \\
            D A^1_i(\bar y_1,\bar y_2,\bar p_i^j,\bar h) =& - \bar p_{1,t}^i - \frac{\partial}{\partial x}\left(a'(y_{1})\bar y_{1} p^i_{1,x} + a(y_{1}) \bar p^i_{1,x}\right) \\
            &+ \left( a''(y_{1})\bar y_{1} y_{1,x} p^i_{1,x} + a'(y_{1})\bar y_{1}  p^i_{1,x} + a'(y_1)y_{1,x} \bar p^i_{1,x} \right) \\
            &+ D_{11}^2 F_1(y_1,y_2)\bar y_1 p^i_1 + D_{12}^2 F_1(y_1,y_2)\bar y_2 p^i_1 \\
            &+ D_{11}^2 F_2(y_1,y_2)\bar y_1 p^i_2 + D_{12}^2 F_2(y_1,y_2)\bar y_2 p^i_2 \\
            %+ F_1'(y_1)\bar p^i_1 +c_{21}\bar p^i_2 
            &+ D_1 F_1(y_1,y_2) \bar p^i_1 + D_1 F_2(y_1,y_2) \bar p^i_2 - \alpha_i \bar y_1 \chi_{\omega_{i,d}}, \\
		      %D\mathcal{A}^1_i(\bar y_1,\bar y_2,\bar p_i^j,\bar h) =& - \bar p_{1t}^i - (a(x) \bar p_{1x}^i)_x +\left(B\sqrt{a} \bar p_1^i\right)_x+ F_1''(y_1)\bar y_1 p^i_1 + F_1'(y_1)\bar p^i_1 +c_{21}\bar p^i_2- \alpha_i \bar y_1 \cara_{\mathcal{O}_{id}}, \\
		      D A^2_i(\bar y_1,\bar y_2,\bar p_i^j,\bar h) =& - \bar p_{2t}^i - \frac{\partial}{\partial x}\left(a'(y_{2})\bar y_{2} p^i_{2,x} + a(y_{2}) \bar p^i_{2,x}\right) \\
            &+ \left( a''(y_{2})\bar y_{2} y_{2,x} p^i_{2,x} +  a'(y_{2})\bar y_{2,x}  p^i_{2,x} + a'(y_2)y_{2,x} \bar p^i_{2,x} \right)\\
            &+ D_{21}^2 F_1(y_1,y_2)\bar y_1 p^i_1 + D_{22}^2 F_1(y_1,y_2)\bar y_2 p^i_1 \\
            &+ D_{21}^2 F_2(y_1,y_2)\bar y_1 p^i_2 + D_{22}^2 F_2(y_1,y_2)\bar y_2 p^i_2 \\
            %+ F_1'(y_1)\bar p^i_1 +c_{21}\bar p^i_2 
            &+ D_2 F_1(y_1,y_2) \bar p^i_1 + D_2 F_2(y_1,y_2) \bar p^i_2 
            %F_2''(y_2)\bar y_2 p^i_2 + F_2'(y_2)\bar p^i_2 
            - \alpha_i \bar y_2 \chi_{\omega_{i,d}}, \\
            %D\mathcal{A}^2_i(\bar y_1,\bar y_2,\bar p_i^j,\bar h) =& - \bar p_{2t}^i - (a(x) \bar p_{2x}^i)_x +\left(B\sqrt{a} \bar p_2^i\right)_x+ F_2''(y_2)\bar y_2 p^i_2 + F_2'(y_2)\bar p^i_2 - \alpha_i \bar y_2 \cara_{\mathcal{O}_{id}}, \\
		      D A^1_3(\bar y_1,\bar y_2,\bar p_i^j,\bar h) =& \, \bar y_1(0),\\
		      D A^2_3(\bar y_1,\bar y_2,\bar p_i^j,\bar h) =& \, \bar y_2(0).
	    %\end{split}\right.    
        \end{cases}
    \end{equation*}We have to show that, for $i=0,1,2,3$ and $j=1,2$, 
    $$
        \frac{1}{\lambda}\left[ A^j_i (( y_1, y_2, p_i^j, h)+\lambda(\bar y_1,\bar y_2,\bar p_i^j,\bar h)) - A^j_i ( y_1, y_2, p_i^j, h) \right] \to D A^j_i(\bar y_1,\bar y_2,\bar p_i^j,\bar h),
    $$
    strongly in the corresponding factor of $Z$ as $\lambda \to 0$.

    Initially, we have
    \begin{equation*}		
    	\begin{split}
    		&\left\| \frac{1}{\lambda}\left[ A^1_0 (( y_1, y_2, p_i^j, h)+\lambda(\bar y_1,\bar y_2,\bar p_i^j,\bar h)) - A^1_0( y_1, y_2, p_i^j, h) \right] - D A^1_0(\bar y_1,\bar y_2,\bar p_i^j,\bar h) \right\|^{2}_{\mathcal{M}} \\
    		&=\left\|
    		\frac{1}{\lambda}\Big[ y_{1,t}+\lambda \bar y_{1,t} - \left( a(y_1 + \lambda \bar y_1) (y_{1,x}+\lambda \bar y_{1,x}) \right)_x + F_1(y_1+\lambda \bar y_1, y_2+\lambda \bar y_2)- (h+\lambda \bar h) \chi_{\omega} \right.   \\
    		& \left. \left.    + \frac{1}{\mu_1} \rho_*^{-2} (p^1_1+\lambda \bar p^1_1) \chi_{\omega_1} + \frac{1}{\mu_2} \rho_*^{-2} (p^2_1+\lambda \bar p^2_1) \chi_{\omega_2} -y_{1,t} + \left( a(y_1) y_{1,x} \right)_x - F_1(y_1,y_2) + h \chi_{\omega}   \right. \right.\\
    		& \left. \left.    - \frac{1}{\mu_1} \rho_*^{-2} p^1_1 \chi_{\omega_1} - \frac{1}{\mu_2} \rho_*^{-2} p^2_1 \chi_{\omega_2} \right] -\bar y_{1,t} + (a'(y_{1})\bar y_{1} y_{1,x} +a(y_1) \bar y_{1,x})_x - D_1 F_1(y_1,y_2)\bar y_1 \right.\\
    		&  \left.  - D_2 F_1(y_1,y_2)\bar y_2 + \bar h \chi_{\omega} \right.  \left. -\frac{1}{\mu_1} \rho_*^{-2} \bar p^1_1 \chi_{\omega_{1}} -\frac{1}{\mu_2} \rho_*^{-2} \bar p^2_1 \chi_{\omega_{2}} \right\|^{2}_{\mathcal{M}} \\
    		& = \int_Q \tilde\rho^2 \left| \frac{1}{\lambda}\left( - \left( a(y_1 + \lambda \bar y_1) (y_{1,x}+\lambda \bar y_{1,x}) \right)_x + \left( a(y_1)y_{1,x} \right)_x \right) + (a'(y_{1})\bar y_{1} y_{1,x} +a(y_1) \bar y_{1,x})_x \right|^2 \\
            & + \int_Q \tilde\rho^2 \left| \frac{1}{\lambda}( F_1(y_1+\lambda\bar y_1,y_2 + \lambda\bar y_2) - F_1(y_1,y_2)) - D_1 F_1(y_1,y_2)\bar y_1 - D_2 F_1(y_1,y_2)\bar y_2 \right|^2 \\
            &= J_1+J_2.
    		%& + \int_Q \rho_2^2 \left| l(\int_0^1 (y+\lambda\bar y)) - l(\int_0^1 y) \right|^2 |(a(x)\bar y_x)_x|^2 = J_1 + J_2,
    	\end{split}
    \end{equation*}From the Mean Value Theorem, there exists $\theta \in [0,1]$, such that
    \begin{equation*}
        \begin{split}
            J_1 \leq &\int_Q \tilde\rho^2 \left( \left| a''(y_1 + \theta \lambda \bar y_1) - a''(y_1) \right|^2 |\bar y_1 y_{1,x} |^2 - \left| a'(y_1 + \lambda \bar y_1)-a'(y_1) \right|^2 \left| 2 y_{1,x} \bar y_{1,x} \right|^2 \right.\\
            &\left.  + \left| a'(y_1+\lambda\bar y_1)-a'(y_1) \right|^2 |\bar y_1 y_{1,xx}|^2 + |a(y_1+\lambda\bar y_1)-a(y_1)|^2 |\bar y_{1,xx}|^2 \right) 
        \end{split}
    \end{equation*}Since $a \in C^3(I)$,  from estimate abive and Lebesgue's dominated convergence theorem, we get that $J_1$ converges to zero as $\lambda \to 0$.
    
    \noindent Once again, from the Mean Value Theorem, there exists 
    %$c \in [a,b]$, where $a=\min \{y_1, y_1 +\lambda \bar y_1 \}$ and $b(t) = \max \{ y_1 , y_1 +\lambda \bar y_1 \}$, 
    $\theta \in [0,1]$, such that
    \begin{eqnarray*}
    %	J_1&=&\int_Q \rho_2^2 \left| \frac{1}{\lambda}(F_1(y_1+\lambda\bar y_1)-F_1(y_1)) - F_1'(y_1)\bar y_1 \right|^2\\
    	J_2&=&\int_Q \tilde\rho^2 \left| \frac{1}{\lambda} \nabla F_1(y_1+\theta\lambda\bar y_1,y_2+\theta\lambda\bar y_2)\cdot (\lambda \bar y_1,\lambda \bar y_2) - D_1 F_1(y_1,y_2)\bar y_1 - D_2 F_1(y_1,y_2)\bar y_2 \right|^2\\
        &=& \int_Q \tilde\rho^2 \left| \nabla F_1(y_1+\theta\lambda\bar y_1,y_2+\theta\lambda\bar y_2)\cdot (\bar y_1,\bar y_2) - \nabla F_1(y_1,y_2) \cdot (\bar y_1,\bar y_2) \right|^2\\
        &=& \int_Q \tilde\rho^2 \left| \nabla F_1(y_1+\theta\lambda\bar y_1,y_2+\theta\lambda\bar y_2) - \nabla F_1(y_1,y_2) \right|^2 (\bar y_1^2 + \bar y_2^2 ).
    	%&=&\int_Q \rho_2^2 \left| F_1'(c) \bar y_1 - F'_1(y_1)\bar y_1 \right|^2.
    \end{eqnarray*}
    Therefore, from \eqref{condiciones_a4f}, we obtain that $J_2$ converges to zero as $\lambda \to 0$. 
    
\noindent On the other hand, for $D A^2_0(\bar y_1,\bar y_2,\bar p_i^j,\bar h)$, we get
    \begin{equation*}		
    	\begin{split}
    		&\left\| \frac{1}{\lambda}\left[ A^2_0 (( y_1, y_2, p_i^j, h)+\lambda(\bar y_1,\bar y_2,\bar p_i^j,\bar h)) - A^2_0( y_1, y_2, p_i^j, h) \right] - D A^2_0(\bar y_1,\bar y_2,\bar p_i^j,\bar h) \right\|^{2}_{\mathcal{M}} \\
    		& = \int_Q \tilde\rho^2 \left| \frac{1}{\lambda}\left( - \left( a(y_2 + \lambda \bar y_2) (y_{2,x}+\lambda \bar y_{2,x}) \right)_x + \left( a(y_2)y_{2,x} \right)_x \right) + (a'(y_{2})\bar y_{2} y_{2,x} +a(y_2) \bar y_{2,x})_x \right|^2 \\ 
            &\quad + \int_Q \tilde\rho^2 \left| \frac{1}{\lambda}(F_2(y_1+\lambda\bar y_1,y_2+\lambda\bar y_2)-F_2(y_1,y_2)) - D_1 F_2(y_1,y_2)\bar y_1 - D_2 F_2(y_1,y_2)\bar y_2 \right|^2 \\
            & = J_3+J_4.
    		%& + \int_Q \rho_2^2 \left| l(\int_0^1 (y+\lambda\bar y)) - l(\int_0^1 y) \right|^2 |(a(x)\bar y_x)_x|^2 = J_1 + J_2,
    	\end{split}
    \end{equation*}As above, from the Mean Value Theorem and Lebesgue's dominated convergence theorem, it follows that
    \begin{equation*}
        \begin{split}
            J_3 \leq &\int_Q \tilde\rho^2 \left( \left| a''(y_2 + \theta \lambda \bar y_2) - a''(y_2) \right|^2 |\bar y_2 y_{2,x} |^2 - \left| a'(y_2 + \lambda \bar y_2)-a'(y_2) \right|^2 \left| 2 y_{2,x} \bar y_{2,x} \right|^2 \right.\\
            &\left. \qquad + \left| a'(y_2 + \lambda\bar y_2) - a'(y_2) \right|^2 |\bar y_2 y_{2,xx}|^2 + |a(y_2+\lambda\bar y_2)-a(y_2)|^2 |\bar y_{2,xx}|^2 \right) \to 0,
        \end{split}
    \end{equation*} as $\lambda \to 0$ and some  constant $\theta \in [0,1]$ 
    
   \noindent Analogously, %$J_2=\int_Q \rho_2^2 \left| F_2'(c) \bar y_2 - F_2'(y_2)\bar y_2 \right|^2$,
    $$
        J_4 = \int_Q \tilde\rho^2 \left| \nabla F_1(y_1+\theta\lambda\bar y_1,y_2+\theta\lambda\bar y_2) - \nabla F_1(y_1,y_2) \right|^2 (\bar y_1^2 + \bar y_2^2 ) \to 0,
    $$ as $\lambda \to 0$.
    
    \noindent Moreover, for $A^1_i$, $i=1,2$, we get
    \begin{equation}
    	\begin{split}
    		&\left\| \frac{1}{\lambda}\left[ A^1_i (( y_1, y_2, p_i^j, h)+\lambda(\bar y_1,\bar y_2,\bar p_i^j,\bar h)) - A^1_i( y_1, y_2, p_i^j, h) \right] - D A^1_i(\bar y_1,\bar y_2,\bar p_i^j,\bar h) \right\|^{2}_{\mathcal{M}} \\
            &= \left\| \frac{1}{\lambda} \left[ 
            - \left(a(y_1+\lambda\bar y_1)(p_{1,x}^1 + \lambda \bar p_{1,x}^1) + a'(y_1+\lambda\bar y_1)(y_{1,x}+\lambda\bar y_{1,x})(p_{1,x}^1+\lambda\bar p_{1,x}^1) \right.\right.\right. \\
            &\left.\left.\left. \qquad + a(y_1)p_{1,x}^1 - a'(y_1)y_{1,x}p_{1,x}^1 \right)_x \right] + \left(a'(y_{1})\bar y_{1} p^i_{1,x} + a(y_{1}) \bar p^i_{1,x}\right)_x \right.\\
            &\qquad \left. - \left( a''(y_{1})\bar y_{1} y_{1,x} p^i_{1,x} + a'(y_{1})\bar y_{1}  p^i_{1,x} + a'(y_1)y_{1,x} \bar p^i_{1,x} \right)
            \right\|_{L^2(\tilde\rho^2,Q)}^2 \\
            &\quad + \left\| \frac{1}{\lambda} \left[ D_1 F_1(y_1+\lambda \bar y_1,y_2+\lambda \bar y_2) (p^i_1 +\lambda \bar p^i_1) + D_1 F_2(y_1+\lambda \bar y_1,y_2+\lambda \bar y_2) (p^i_2 +\lambda \bar p^i_2) \right.\right. \\
            &\left.\left.\qquad - D_1 F_1(y_1,y_2) p^i_1 - D_1 F_2(y_1,y_2) p^i_2 \right] - D_{11}^2 F_1(y_1,y_2)\bar y_1 p^i_1 - D_{12}^2 F_1(y_1,y_2)\bar y_2 p^i_1 \right. \\
            &\left.\qquad - D_{11}^2 F_2(y_1,y_2)\bar y_1 p^i_2 - D_{12}^2 F_2(y_1,y_2)\bar y_2 p^i_2 
            - D_1 F_1(y_1,y_2) \bar p^i_1 - D_1 F_2(y_1,y_2) \bar p^i_2
            %- F''_1(y_1)\bar y_1 p^i_1 - F_1'(y_1)\bar p^i_1
            \right\|^{2}_{\mathcal{M}} \\
            &= J_5+J_6.
%    		&= \left\| \frac{1}{\lambda} \left( F_1'(y_1+\lambda \bar y_1) (p^i_1 +\lambda \bar p^i_1) - F_1'(y_1) p^i_1  \right) - F''_1(y_1)\bar y_1 p^i_1 - F_1'(y_1)\bar p^i_1 \right\|^{2}_{L^2(\rho_2^2,Q)} = J_5+J_6. 
    	\end{split}
    \end{equation}We observe that  $J_5$  is formed by two terms: $J_{51}$ and $J_{51}$, satisfying
    \begin{equation*}
        \begin{split}
            J_{51} \leq 2\left\| \left( \frac{1}{\lambda} \left[ a(y_1+\lambda\bar y_1)(p^1_{1,x}+\lambda\bar p^i_{1,x} ) - a(y_1)p^i_{1,x} \right] - a'(y_1)\bar y_1 p^i_{1,x} + a(y_1)\bar p^i_{1,x} \right)_x\right\|^{2}_{\mathcal{M}}
        \end{split}
    \end{equation*}
    and
    \begin{equation*}
        \begin{split}
            J_{52} \leq &2\left\| \left( \frac{1}{\lambda} \left[ a'(y_1+\lambda\bar y_1)(y_{1,x}-\lambda\bar y_{1,x})(p^1_{1,x}+\lambda\bar p^i_{1,x} ) - a'(y_1)y_{1,x}p^i_{1,x} \right] \right.\right. \\
            &\left.\left. \qquad - (a''(y_1)\bar y_1 y_{1,x} p^1_{1,x} + a'(y_1)\bar y_{1,x} p^1_{1,x} + a'(y_1)y_{1,x}\bar p^i_{1,x}) \right)_x\right\|^{2}_{\mathcal{M}}.
        \end{split}
    \end{equation*}From the above estimates, calculations similar to the previous ones allow us to obtain that $J_{51}$ and $J_{52}$ converge to zero as $\lambda \to 0$.

    \noindent Since 
    \begin{equation*}
    J_6 \leq     J_{61} + J_{62} + J_{63},
    \end{equation*}where
    \begin{equation*}
    \begin{cases}
    J_{61} &= \displaystyle \int_Q \tilde\rho^2 \Big| \frac{1}{\lambda}\left[D_1 F_1(y_1+\lambda \bar y_1,y_2+\lambda \bar y_2) p^i_1 - D_1 F_1(y_1,y_2) p^i_1 \right] - \\
   & \quad D_{11}^2 F_1(y_1,y_2)\bar y_1 p^i_1 - D_{12}^2 F_1(y_1,y_2)\bar y_2 p^i_1\Big|^2, \\
   J_{62} &= \displaystyle \int_Q \tilde\rho^2 \Big| \frac{1}{\lambda}\left[ D_1 F_2(y_1+\lambda \bar y_1,y_2+\lambda \bar y_2) p^i_2  - D_1 F_2(y_1,y_2) p^i_2 \right] \\
   & \quad - D_{11}^2 F_2(y_1,y_2)\bar y_1 p^i_2 - D_{12}^2 F_2(y_1,y_2)\bar y_2 p^i_2 \Big|^2,\\
   J_{63} &= \displaystyle \int_Q \tilde\rho^2 \Big| D_1 F_1(y_1+\lambda \bar y_1,y_2+\lambda \bar y_2) \bar p^i_1 + D_1 F_2(y_1+\lambda \bar y_1,y_2+\lambda \bar y_2) \bar p^i_2 \\
   & \quad D_1 F_1(y_1,y_2) \bar p^i_1 - D_1 F_2(y_1,y_2) \bar p^i_2 \Big|^2,
    \end{cases}
    \end{equation*}
    
   % On the other hand, 
    %\begin{equation*}
    %	\begin{split}
    %		&J_6\\ %&= \left\| \frac{1}{\lambda} \left( F_1'(y_1+\lambda \bar y_1) (p^i_1 +\lambda \bar p^i_1) - F_1'(y_1) p^i_1  \right) - F''_1(y_1)\bar y_1 p^i_1 - F_1'(y_1)\bar p^i_1 \right\|^{2}_{L^2(\rho_2^2,Q)} \\
            %\int_Q \rho_2^2 \left| \frac{1}{\lambda} \left( F_1'(y_1+\lambda \bar y_1) (p^i_1 +\lambda \bar p^i_1) - F_1'(y_1) p^i_1  \right) - F''_1(y_1)\bar y_1 p^i_1 - F_1'(y_1)\bar p^i_1 \right|^2 \\
    		%&\leq \int_Q \tilde\rho^2 \left| \frac{1}{\lambda}\left[D_1 F_1(y_1+\lambda \bar y_1,y_2+\lambda \bar y_2) p^i_1 - D_1 F_1(y_1,y_2) p^i_1 \right] - D_{11}^2 F_1(y_1,y_2)\bar y_1 p^i_1 - D_{12}^2 F_1(y_1,y_2)\bar y_2 p^i_1\right|^2  \\
            %&+ \int_Q \tilde\rho^2 \left| \frac{1}{\lambda}\left[ D_1 F_2(y_1+\lambda \bar y_1,y_2+\lambda \bar y_2) p^i_2  - D_1 F_2(y_1,y_2) p^i_2 \right] - D_{11}^2 F_2(y_1,y_2)\bar y_1 p^i_2 - D_{12}^2 F_2(y_1,y_2)\bar y_2 p^i_2 \right|^2 \\
            %&+\int_Q \tilde\rho^2 \left| D_1 F_1(y_1+\lambda \bar y_1,y_2+\lambda \bar y_2) \bar p^i_1 + D_1 F_2(y_1+\lambda \bar y_1,y_2+\lambda \bar y_2) \bar p^i_2  - D_1 F_1(y_1,y_2) \bar p^i_1 - D_1 F_2(y_1,y_2) \bar p^i_2 \right|^2 \\
            %&= J_{61} + J_{62} + J_{63}.
        %\end{split}
   % \end{equation*}
    \noindent from the Mean Value Theorem, for $\theta \in [0,1]$,
    \begin{equation*}
        \begin{split}
            J_{61} &\leq \int_Q \tilde\rho^2 \left| \frac{1}{\lambda}
            \nabla (D_1 F_1) (y_1+ \theta\lambda \bar y_1,y_2+\theta\lambda \bar y_2)\cdot(\lambda \bar y_1,\lambda \bar y_2) p^i_1 - \nabla (D_1 F_1)(y_1,y_2)\cdot(\bar y_1,\bar y_2)p^i_1
            \right|^2,
        \end{split}
    \end{equation*}
    and the fact that $D_1 F \in C^1$, we get that $J_{31}$ converges to zero as $\lambda \to 0$. 
    
   \noindent  From similar computations, we find  that $J_{62}$ and $F_{63}$ converges to zero as $\lambda \to 0$
    
    For $A^2_i$, $i=1,2$, a similar estimation gives
    $$
        \left\| \frac{1}{\lambda}\left[ A^1_i (( y_1, y_2, p_i^j, h)+\lambda(\bar y_1,\bar y_2,\bar p_i^j,\bar h)) - A^1_i( y_1, y_2, p_i^j, h) \right] - D A^1_i(\bar y_1,\bar y_2,\bar p_i^j,\bar h) \right\|^{2}_{\mathcal{M}} \to 0,
    $$
    as $\lambda \to 0$. This finishes the proof that $A$ is Gateaux differentiable, with a \textit{G-derivative} given by $D A( y_1, y_2, p_i^j, h)$.
    
     Now, take $( y_1, y_2, p_i^j, h)\in Y$ and let $\{(y_{1,n},y_{2,n},p^{j}_{i,n},h_{n})\}_{n=0}^{\infty}$ be a sequence that converges to  $( y_1, y_2, p_i^j, h)$ in $Y$. For each $(\bar y_1,\bar y_2,\bar p_i^j,\bar h)\in B_{r}(0)$, we have to prove that 
    \begin{equation}\label{eq:continuity_DA_m_l}
    	\begin{split}
    		&\left\|(D A^l_m(y_{1,n},y_{2,n},p^{j}_{i,n},h_{n})-D A^l_m( y_1, y_2, p_i^j, h))(\bar y_1,\bar y_2,\bar p_i^j,\bar h) \right\|^{2}_{\mathcal{M}}\rightarrow 0,
    	\end{split}
    \end{equation}
    as $\lambda \to 0$, for $m=0,1,2,3$ and $l=1,2$.
    
    We present the details for $m=0$ and $l=1$. The other terms are estimated analogously, and we omit them. We have
    \begin{equation*}
        \begin{split}
    	   &(D A^1_0(y_{1,n},y_{2,n},p^{j}_{i,n},h_{n})-D A^1_0( y_1, y_2, p_i^j, h))(\bar y_1,\bar y_2,\bar p_i^j,\bar h)\\
           &= \left(a'(y_{1,n})\bar y_{1} y_{1,nx} - a'(y_{1})\bar y_{1} y_{1,x} + a(y_{1,n}) \bar y_{1,x} - a(y_1) \bar y_{1,x} \right)_x \\
           &\quad + [ (D_1 F_1(y_{1,n},y_{2,n})-D_1 F_1(y_1,y_2))\bar y_1 + (D_2 F_1(y_{1,n},y_{2,n})-D_2 F_1(y_1,y_2))\bar y_2] \\
           &= ((a'(y_{1,n})-a'(y_{1}))\bar y_{1} y_{1,nx})_x + (a'(y_{1,n})\bar y_{1} (y_{1,nx}-y_{1,x}))_x + ((a(y_{1,n})-a(y_{1})) \bar y_{1,x})_x \\&\quad + [ (D_1 F_1(y_{1,n},y_{2,n})-D_1 F_1(y_1,y_2))\bar y_1 + (D_2 F_1(y_{1,n},y_{2,n})-D_2 F_1(y_1,y_2))\bar y_2] \\
           &= X_0^{11}+X_0^{12}+X_0^{13}+X_0^{14}.
        \end{split}
    \end{equation*}First, 
    \begin{eqnarray*}
        \int_{Q} \tilde\rho^{2}|X_{0}^{11}|^{2} dxdt
        %&\leq&  
        %\int_{Q}\tilde\rho^{2} \left|((a'(y_{1,n})-a'(y_{1}))\bar y_{1} y_{1,nx})_x\right|^2 + \left|(a'(y_{1,n})\bar y_{1} (y_{1,nx}-y_{1,x}))_x \right|^2dxdt\\
        %&& + \iint_{Q}\tilde\rho^{2} \left|((a(y_{1,n})-a(y_{1})) \bar y_{1,x})_x \right|^2 dxdt\\
    	&\leq&C\int_{Q} \tilde\rho^{2} M^2 \left( \left|(y_{1,n}-y_1)_x \bar y_1 y_{1,nx} \right|^2 + \left|(y_{1,n}-y_1) \bar y_{1,x} y_{1,nx}  \right|^2 \right. \\
        &&\left. \qquad\quad +\left|(y_{1,n}-y_1) \bar y_1 y_{1,nxx} \right|^2 \right) dxdt.
    \end{eqnarray*}
    Using embedding $H^1_0(I) \hookrightarrow L^\infty(I)$ and %that $\tilde \rho^2 \hat\rho_1^{-6} \leq C$, see 
    \eqref{eq:compara_rhos3}, we get
    \begin{equation*}
        \begin{split}
            &\int_{Q} \tilde\rho^{2} \left|(y_{1,n}-y_1)_x \bar y_1 y_{1,nx} \right|^2 dx dy \\
            &\leq C\int_0^T \tilde\rho^{2} \|y_{1,nx}-y_{1,x}\|_{L^\infty(I)}^2 \|\bar y_1\|_{L^\infty(I)}^2 \|y_{1,nx}\|_{L^2(I)}^2 dt \\
            &\leq C \int_0^T \tilde\rho^{2} \hat\rho_1^{-6} \hat\rho_1^2 \|y_{1,nxx}-y_{1,xx}\|_{L^2(I)}^2 \hat\rho_1^2 \|\bar y_{1,x}\|_{L^2(I)}^2 \hat\rho_1^2 \|y_{1,nx}\|_{L^2(I)}^2 dt \\
            &\leq C \sup_{t\in[0,T]} \hat\rho_1^2 \|\bar y_{1,x}\|_{L^2(I)}^2 \sup_{t\in[0,T]} \hat\rho_1^2 \|y_{1,nx}\|_{L^2(I)}^2  \int_0^T \hat\rho_1^2 \|y_{1,nxx}-y_{1,xx}\|_{L^2(I)}^2 dt \\
            &\leq C \|(\bar y_1,\bar y_2,\bar p_i^j,\bar h)\|_{Y} \|(y_1,y_2,p_i^j,h)\|_{{Y}} \| (y_{1,n},y_{2,n},p_{i,n}^j,h_n) - (y_1,y_2,p_i^j,h) \|_{{Y}}.
        \end{split}
    \end{equation*}Proceeding similarly with the other terms of $X_0^{11}$, we get
    \begin{equation*}
        \begin{split}
                &\int_{Q} \tilde\rho^{2}|X_{0}^{11}|^{2} dxdt \\
                &\leq C \|(\bar y_1,\bar y_2,\bar p_i^j,\bar h)\|_{{Y}} \|(y_1,y_2,p_i^j,h)\|_{{Y}} \| (y_{1,n},y_{2,n},p_{i,n}^j,h_n) - (y_1,y_2,p_i^j,h) \|_{{Y}}.
        \end{split}
    \end{equation*}Analogously, for $X_0^{12}$ and $X_0^{13}$, 
    \begin{equation*}
        \begin{split}
            &\int_{Q} \tilde\rho^{2} (|X_{0}^{12}|^{2} + |X_{0}^{13}|^{2} )dxdt\\
            &\leq C\int_{Q} \tilde\rho^{2} M^2 \left( |(\bar y_1 (y_{1,nx}-y_{1,x}))_x|^2 + |((y_{1,n}-y_1) \bar y_{1,x})_x|^{2} \right) dxdt \\
            &\leq C\|(\bar y_1,\bar y_2,\bar p_i^j,\bar h)\|_{{Y}} \| (y_{1,n},y_{2,n},p_{i,n}^j,h_n) - (y_1,y_2,p_i^j,h) \|_{{Y}}.
        \end{split}
    \end{equation*}
    Therefore,
    \begin{equation*}
        \int_{Q} \tilde\rho^{2} (|X_{0}^{11}|^{2} + |X_{0}^{12}|^{2} + |X_{0}^{13}|^{2} )dxdt \to 0, \qquad\text{as}\quad n\to \infty.
    \end{equation*}
    On the other hand,
    \begin{equation*}
        \begin{split}
        	\int_{Q} \tilde\rho^{2} |X_{0}^{14}|^{2} dxdt&\leq  \int_{Q} \tilde\rho^{2} |(D_1 F_1(y_{1,n},y_{2,n})-D_1 F_1(y_1,y_2))|^2 |\bar y_1|^{2}dxdt\\
            & + \int_{Q} \tilde\rho^{2} |(D_2 F_1(y_{1,n},y_{2,n})-D_2 F_1(y_1,y_2))|^2 |\bar y_2|^{2}dxdt\\
        	&\leq C\int_{Q} \tilde\rho^{2} M^2 ( |y_{1,n}-y_1|^2 + |y_{2,n}-y_2|^2 ) (|\bar y_1|^{2} + |\bar y_2|^{2} ) dxdt.
        \end{split}
    \end{equation*}
    %\color{red}
    
    \noindent From  \eqref{eq:compara_rhos3}, we have  that  $\tilde\rho^2\leq C \rho_1^2 \rho^2_1\leq C \hat{\rho}^2 \rho_0^2 $. Then
    \begin{equation*}
    \begin{split}
    	&\int_{Q} \tilde\rho^{2}|X_{0}^{14}|^{2} dxdt\\
        &\leq C \int_{Q} \hat{\rho}_1^2 \rho_0^2  ( |y_{1,n}-y_1|^2 + |y_{2,n}-y_2|^2 ) (|\bar y_1|^{2} + |\bar y_2|^{2} )dxdt\\
    	&\leq C \int_0^T \left(\hat{\rho}_1^2 (\| \bar{y}_1\|^2_{L^\infty(I)} + \| \bar{y}_2\|^2_{L^\infty(I)})\right)\int_{0}^{1} \rho_0^2  (|y_{1,n}-y_1|^2 + |y_{2,n}-y_2|^2 ) dt\\
        &\leq C \int_0^T \sup_{t\in [0,T]}\left(\hat{\rho}_1^2 (\| \bar{y}_{1,x}\|^2_{L^2(I)} + \| \bar{y}_{2,x}\|^2_{L^2(I)})\right) \left( \|\rho_0 (y_{1,n}-y_1)\|_{L^2(I)} + \|\rho_0 (y_{2,n}-y_2)\|_{L^2(I)} \right) dt\\
    	&\leq C\|(\bar y_1,\bar y_2,\bar p_i^j,\bar h)\|_{{Y}} \cdot \left( \|\rho_0 (y_{1,n}-y_1)\|_{L^2(Q)} + \|\rho_0 (y_{2,n}-y_2)\|_{L^2(Q)} \right)\\
    	&\leq C\|(y_{1,n}-y_1),(y_{2,n}-y_2),(p^{j}_{i,n}-p^{j}_i),(h_{n}-h)\|_{{Y}}\rightarrow 0 \quad\text{as}\quad n \to \infty.
    \end{split}
    \end{equation*}
    This shows that \eqref{eq:continuity_DA_m_l} is satisfied. 
\end{proof}

\end{document}
